
En este capítulo se presentan los objetivos y el catálogo de requisitos del nuevo sistema informático. Opcionalmente, se incluye el análisis de la brecha entre los requisitos planteados y la solución base seleccionada. 

\section{Objetivos de Negocio}
\textcolor{red}{TODO}Esta sección debe contener los objetivos de negocio que se esperan alcanzar cuando el sistema software a desarrollar esté en producción. Adicionalmente, se pueden incluir los modelos de procesos, que normalmente son los modelos de procesos de negocio actuales con ciertas mejoras.

\section{Objetivos del Sistema}
\textcolor{red}{TODO} Esta sección debe contener la especificación de los objetivos o requisitos generales del sistema.

\section{Metodología de captura de requisitos}
\label{sec:RequirementsCaptureMethodology}
\subsection{Reunión de toma de requisitos}

Con el objetivo de identificar de forma precisa las necesidades del cliente y establecer una base sólida para el análisis y diseño del sistema, se llevó a cabo una reunión inicial de toma de requisitos el día 03/03/2026.

La obtención de requisitos se realizó mediante una entrevista con el interesado principal del sistema, una técnica habitual dentro del proceso de \textit{requirements elicitation} descrito en el estándar IEEE 29148 \citep{ieee29148}. En este caso se utilizó una entrevista semiestructurada, apoyada en un guion previo que permitió mantener cierta flexibilidad durante la conversación y profundizar en aquellos aspectos que surgían durante la sesión \citep[pp.~104--106]{sommerville}.

Durante la entrevista se abordaron de manera sistemática distintos aspectos relacionados con el contexto del negocio, los procesos actuales y las necesidades del sistema, entre los que destacan:

\begin{itemize}
    \item Contexto del negocio, objetivos de la aplicación y problemática actual en la gestión comercial.
    \item Relación operativa con la empresa distribuidora central y flujo actual de pedidos.
    \item Tipos de usuarios implicados en el sistema (comercial, peluquerías/clientes profesionales y clientes finales).
    \item Gestión del catálogo de productos capilares y condiciones comerciales asociadas.
    \item Proceso de venta, pedidos, descuentos y seguimiento de clientes.
    \item Necesidades de gestión comercial, seguimiento de tratamientos y fidelización de clientes.
\end{itemize}

Con el fin de garantizar la fidelidad de la información recogida, la entrevista fue grabada en audio mediante un dispositivo Plaud Note \citep{plaudnote}, que permitió generar una transcripción automática posterior. Dicha transcripción fue revisada manualmente para asegurar la correcta interpretación de los requisitos expresados.

La información obtenida durante esta sesión ha servido como base para la elaboración de los modelos conceptuales, la definición de casos de uso y la específicación de los requisitos funcionales y no funcionales que se detallan en los apartados siguientes.

\subsection{Análisis de tecnologías existentes}

Como parte del proceso de análisis del sistema, se ha realizado una revisión de las herramientas y aplicaciones utilizadas actualmente por el distribuidor en su actividad diaria. Se han extraído gran parte de los campos de información necesarios para el diseño del nuevo sistema, facilitando así la definición del modelo de datos y garantizando la viabilidad de la migración de la información existente.

Entre las herramientas analizadas se encuentran aplicaciones de gestión de bases de datos y hojas de cálculo, como \textit{Microsoft Access} y \textit{Microsoft Excel}, así como el software de facturación \textit{Factusol}, ampliamente utilizado en la operativa actual del distribuidor.

El estudio de estas tecnologías ha permitido identificar tanto buenas prácticas reutilizables como carencias funcionales, que han sido tenidas en cuenta en la definición de los requisitos del nuevo sistema.

\section{División modular de requisitos}
\label{sec:ModularRequirements}
Con el objetivo de facilitar la comprension del sistema y organizar las funcionalidades de forma coherente, los requisitos se han agrupado en modulos funcionales que representan las principales areas de operacion del sistema. Estos modulos se describen a continuacion de forma general, mientras que en los apartados posteriores se detallan sus requisitos funcionales y de informacion asociados.

La descomposicion del sistema en modulos funcionales responde a los principios de diseno arquitectonico en ingenieria del software, que promueven la organizacion del sistema en componentes con responsabilidades bien definidas, alta cohesion y bajo acoplamiento, con el objetivo de facilitar su comprension, mantenimiento y evolucion \citep[Ch.~6--7]{sommerville}.

\subsection{M1. Gestion de acceso, alta y autorizacion de usuarios}

Este modulo agrupa las funcionalidades relacionadas con la identificacion de los usuarios del sistema, el alta de nuevas cuentas, la autenticacion y el control de acceso segun el perfil asignado. Su finalidad es garantizar que unicamente usuarios autorizados puedan utilizar la aplicacion y que cada uno acceda exclusivamente a las funciones que le correspondan dentro del sistema.

Asimismo, este modulo contempla la gestion de roles, el tratamiento de solicitudes de registro, la administracion del estado de las cuentas y el registro de eventos de acceso y acciones administrativas, con el fin de aportar seguridad, trazabilidad y control sobre el uso de la plataforma.

Este modulo constituye la base del sistema, ya que establece los mecanismos de acceso y control sobre los que se apoyan el resto de funcionalidades. Su correcta definicion permite garantizar la seguridad del sistema, la coherencia en la gestion de usuarios y la trazabilidad de las operaciones realizadas.

\subsection{M2. Gestion comercial y clientes}

Este modulo agrupa las funcionalidades relacionadas con la gestion de los clientes profesionales (peluquerias) y con la operativa diaria del equipo comercial. Su objetivo es centralizar la informacion relevante de cada cliente, mantener controlada su asignacion a un comercial concreto y facilitar la planificacion del trabajo de campo mediante visitas, rutas y estimaciones operativas de reparto.

Ademas de la informacion basica y comercial de cada cliente, este modulo incorpora la geolocalizacion de los establecimientos, la definicion de franjas horarias validas para visitas, la configuracion operativa de la jornada de cada comercial y la generacion dinamica de una planificacion diaria. De este modo, el sistema no actua unicamente como repositorio de fichas de cliente, sino tambien como herramienta de apoyo a la toma de decisiones durante la jornada, calculando orden de visitas, horas estimadas de llegada, margenes disponibles y situaciones que requieren replanificacion.

\subsection{M3. Catalogo, productos y coloracion}

Este modulo agrupa las funcionalidades relacionadas con la consulta y gestion del catalogo de productos profesionales comercializados por el distribuidor. Su objetivo es proporcionar a los clientes profesionales y a los agentes comerciales una forma estructurada de acceder a la informacion de los productos, sus caracteristicas tecnicas y su clasificacion dentro de las distintas lineas comerciales.

Asimismo, el modulo contempla la organizacion jerarquica del catalogo mediante categorias y lineas comerciales, asi como la gestion de recursos de apoyo, documentacion tecnica y materiales complementarios asociados a productos o gamas. Tambien incluye herramientas especificas para la consulta de cartas de color y referencias cromaticas utilizadas en los procesos de coloracion capilar, lo que refuerza su utilidad tanto comercial como tecnica.

\subsection{M4. Pedidos, entregas y cobros}

Este modulo agrupa las funcionalidades relacionadas con la gestion del proceso comercial de venta de productos, desde la creacion de borradores de pedido hasta la entrega fisica y el seguimiento minimo del cobro asociado. Su objetivo es permitir que clientes profesionales, comerciales y administradores trabajen sobre un mismo flujo trazable de pedido, con visibilidad del estado del encargo, de su reparto y de su situacion de cobro.

En la implementacion actual, el modulo se apoya en pedidos con lineas de producto, estados de ciclo de vida, vinculacion opcional a visitas de reparto y un seguimiento minimo del cobro mediante estados \texttt{pending}/\texttt{paid}, metodo, fecha y notas. La fase logistica no se modela como un subsistema aislado de entregas, sino integrada en las visitas comerciales de tipo reparto, con validacion de cierre mediante QR y con ETA visible para cliente solo cuando existe un reparto real y coherente. Quedan explicitamente fuera de este alcance minimo los pagos parciales, la facturacion y la conciliacion avanzada.

\subsection{M5. Gestion tecnica del salon y seguimiento de servicios}

Este modulo agrupa las funcionalidades destinadas a registrar y consultar informacion tecnica sobre los servicios realizados en las peluquerias clientes del sistema. Su objetivo es permitir a los profesionales del salon llevar un seguimiento estructurado de los trabajos realizados a sus clientes, especialmente aquellos relacionados con productos de la marca, facilitando la consulta de historiales tecnicos y el analisis del uso de productos.

En este sentido, el modulo contempla la gestion de clientes finales del salon, los servicios o tratamientos realizados, sus fichas tecnicas asociadas y los productos empleados en cada intervencion. Ademas, incorpora soporte para la generacion de sugerencias de productos y para la elaboracion de comunicaciones personalizadas a partir de la informacion tecnica registrada, reforzando asi tanto la continuidad del trabajo profesional como la fidelizacion del cliente final.

\subsection{M6. Promociones, notificaciones y formaciones}

Este modulo agrupa las funcionalidades relacionadas con la comunicacion entre el distribuidor y las peluquerias usuarias del sistema. Su objetivo es facilitar la difusion de promociones comerciales, la organizacion de actividades formativas y el envio de notificaciones o recordatorios relevantes para los usuarios de la aplicacion.

Ademas, este modulo permite segmentar clientes mediante categorias comerciales, asociar promociones a productos, lineas o grupos de clientes, y gestionar avisos vinculados a eventos relevantes del sistema. De este modo, no solo cubre la comunicacion promocional, sino tambien la comunicacion operativa y la gestion de acciones formativas e inscripciones.

\subsection{M7. Administracion e integraciones}

Este modulo agrupa las funcionalidades destinadas a la administracion general del sistema y a la integracion con herramientas o servicios externos utilizados en la actividad del distribuidor. Su objetivo es permitir la configuracion y mantenimiento de la aplicacion, facilitar el analisis de la actividad comercial y permitir la interoperabilidad con sistemas externos utilizados en la operativa del negocio.

En particular, este modulo contempla la gestion de parametros generales de configuracion, el registro y seguimiento de integraciones externas y las operaciones de exportacion o sincronizacion de informacion. Asimismo, incorpora soporte a la generacion de propuestas de pedido a proveedor a partir de la informacion registrada en la plataforma o introducida manualmente, alineandose con la operativa habitual del distribuidor.

\subsection{M8. Funcionalidades adicionales}

Este modulo recoge un conjunto de funcionalidades identificadas durante el proceso de analisis de requisitos que por ellas solas no forman parte de ningun modulo principal. Estas funcionalidades buscan ampliar el valor de la aplicacion para los clientes mediante el uso de tecnologias avanzadas y nuevas herramientas de apoyo al trabajo del peluquero.

Entre estas funcionalidades se incluyen, de forma potencial, herramientas de recomendacion de tratamientos, simulacion visual de coloracion, reconocimiento de productos mediante camara y gestion de reservas online por parte de clientes finales. Estas funcionalidades pueden llegar a entrar, o no, en el alcance del proyecto dependiendo de la disponibilidad temporal del mismo. Si no llegasen a incluirse en la primera version del sistema, podrian ser consideradas para futuras actualizaciones o mejoras de la aplicacion.


\section{Requisitos funcionales}
\label{sec:FunctionalRequirements}
\section{Requisitos funcionales}
\label{sec:FunctionalRequirements}

A continuación, la Tabla~\ref{tab:functional-requirements-executive} presenta el catálogo de requisitos funcionales en formato ejecutivo. Los requisitos mantienen identificadores únicos para facilitar su trazabilidad, siguiendo las recomendaciones de identificación y seguimiento propuestas por ISO/IEC/IEEE 29148 \cite{ieee29148}. La versión extensa empleada durante el desarrollo se conserva en el repositorio como copia de apoyo.

{\scriptsize
\setlength{\parskip}{0pt}
\setlength{\tabcolsep}{3pt}
\renewcommand{\arraystretch}{1.10}
\newcommand{\RFModuleHeader}[1]{%
\hline
\rowcolor{black!8}
\multicolumn{4}{|>{\raggedright\arraybackslash}p{0.94\linewidth}|}{\textbf{#1}}\\
\hline
}
\tablecaption{Catálogo ejecutivo de requisitos funcionales agrupado por módulos\label{tab:functional-requirements-executive}}
\tablefirsthead{%
\hline
\textbf{ID} & \textbf{Módulo} & \textbf{Actores} & \textbf{Requisito funcional sintetizado} \\
 \hline}
\tablehead{%
\hline
\textbf{ID} & \textbf{Módulo} & \textbf{Actores} & \textbf{Requisito funcional sintetizado} \\
\hline}
\begin{supertabular}{|>{\raggedright\arraybackslash}p{0.085\linewidth}|>{\raggedright\arraybackslash}p{0.205\linewidth}|>{\raggedright\arraybackslash}p{0.18\linewidth}|>{\raggedright\arraybackslash}p{0.43\linewidth}|}
\RFModuleHeader{M1. Acceso, alta y seguridad de cuentas}
\texttt{RF-M1-01} & Acceso y alta & Cliente, administrador & Registrar solicitudes de alta y permitir la creación directa de cuentas, manteniendo revisión administrativa, aprobación, rechazo, activación, desactivación y trazabilidad de la gestión. \\
\hline
\texttt{RF-M1-02} & Acceso y alta & Todos los roles & Autenticar usuarios mediante credenciales seguras, permitir solo cuentas activas y autorizadas, limitar intentos abusivos y registrar accesos correctos y fallidos. \\
\hline
\texttt{RF-M1-03} & Acceso y alta & Administrador & Gestionar roles de administrador, comercial y cliente profesional, restringiendo navegación, pantallas y operaciones según el rol autenticado. \\
\hline
\texttt{RF-M1-04} & Acceso y alta & Administrador & Aplicar condiciones de autorización comercial sobre cuentas profesionales y conservar el estado operativo de cada usuario junto con su historial administrativo. \\
\hline
\texttt{RF-M1-05} & Acceso y alta & Todos los roles & Crear, mantener y cerrar sesiones de usuario, impidiendo la reutilización de sesiones expiradas o revocadas. \\
\hline
\texttt{RF-M1-06} & Acceso y alta & Todos los roles & Permitir consulta y edición del perfil propio, imagen, datos de contacto y cambio de credenciales conforme a las políticas de seguridad. \\
\hline
\RFModuleHeader{M2. Clientes profesionales, comerciales y rutas}
\texttt{RF-M2-01} & Clientes y comerciales & Administrador & Registrar clientes profesionales, actualizar su información operativa, asignarlos o reasignarlos a comerciales y conservar la trazabilidad de esas asignaciones. \\
\hline
\texttt{RF-M2-02} & Clientes y comerciales & Comercial, cliente & Consultar la ficha operativa del cliente con datos de contacto, localización, observaciones, asignación comercial y contexto de planificación. \\
\hline
\texttt{RF-M2-03} & Clientes y comerciales & Comercial, cliente & Planificar, registrar, completar, aplazar o cancelar visitas comerciales, notificando cambios relevantes y respetando la ventana horaria del cliente. \\
\hline
\texttt{RF-M2-04} & Clientes y comerciales & Comercial & Previsualizar y ordenar rutas comerciales diarias con geolocalización, estimación de tiempos y secuencia de visitas. \\
\hline
\texttt{RF-M2-05} & Clientes y comerciales & Comercial, cliente & Configurar la jornada diaria del comercial y mostrar al cliente una estimación de llegada cuando exista una visita de reparto real vinculada. \\
\hline
\RFModuleHeader{M3. Catálogo, productos y coloración}
\texttt{RF-M3-01} & Catálogo & Cliente, comercial & Consultar el catálogo activo de productos con información comercial, técnica e imagen asociada. \\
\hline
\texttt{RF-M3-02} & Catálogo & Cliente, comercial & Filtrar y buscar productos por categoría, línea, subcategoría, referencia, nombre y estado de disponibilidad. \\
\hline
\texttt{RF-M3-03} & Catálogo & Cliente, comercial & Visualizar cartas de color y referencias cromáticas asociadas a productos y líneas comerciales, incluyendo vista por filas de tono principal, vista por cartas y detalle de tono con producto asociado cuando exista. \\
\hline
\texttt{RF-M3-04} & Catálogo & Administrador & Mantener productos, categorías, líneas, subcategorías, recursos de apoyo, cartas de color y referencias desde administración. \\
\hline
\RFModuleHeader{M4. Pedidos, reparto y cobro}
\texttt{RF-M4-01} & Pedidos y reparto & Cliente, comercial & Crear, editar, confirmar, cancelar y consultar pedidos, manteniendo líneas, productos, referencias de color, imágenes de apoyo, descuentos visibles, método de entrega y estados coherentes. \\
\hline
\texttt{RF-M4-02} & Pedidos y reparto & Comercial, cliente & Preparar uno o varios repartos a partir de pedidos confirmados, seleccionando líneas y cantidades, registrando bultos manuales, generando etiqueta logística y vinculando los repartos comerciales a visitas de entrega. \\
\hline
\texttt{RF-M4-03} & Pedidos y reparto & Comercial, administrador & Validar entregas de repartos comerciales mediante QR escaneado o introducción manual y sincronizar el estado logístico con la visita y el pedido cuando todos sus repartos estén completados. \\
\hline
\texttt{RF-M4-04} & Pedidos y reparto & Cliente, comercial, administrador & Permitir elegir entre entrega por comercial o por agencia, aplicando en el segundo caso un cargo configurable y generando una etiqueta diferenciada sin QR de reparto comercial. \\
\hline
\texttt{RF-M4-05} & Pedidos y reparto & Comercial, administrador & Registrar pagos parciales o totales de pedidos entregados, conservando importe, método, fecha, notas, usuario responsable, historial de pagos y saldo pendiente. \\
\hline
\texttt{RF-M4-06} & Pedidos y reparto & Administrador, comercial, cliente & Consultar la trazabilidad cruzada de pedidos, repartos, etiquetas, entregas y cobros desde las vistas correspondientes a cada rol. \\
\hline
\RFModuleHeader{M5. Gestión técnica del salón}
\texttt{RF-M5-01} & Salón & Cliente profesional & Gestionar fichas de clientes finales del salón con datos de contacto, observaciones e historial asociado. \\
\hline
\texttt{RF-M5-02} & Salón & Cliente profesional & Registrar servicios realizados a clientes del salón con fecha, descripción, resultado y relación con la ficha correspondiente. \\
\hline
\texttt{RF-M5-03} & Salón & Cliente profesional & Registrar información técnica de tratamientos, fórmulas, tonalidades, tiempos, productos usados y observaciones profesionales. \\
\hline
\texttt{RF-M5-04} & Salón & Cliente profesional & Consultar el historial técnico de cada cliente del salón, filtrando servicios, tratamientos y resultados previos. \\
\hline
\texttt{RF-M5-05} & Salón & Cliente profesional & Registrar y consultar productos utilizados en servicios para facilitar seguimiento técnico y reposición. \\
\hline
\texttt{RF-M5-06} & Salón & Cliente profesional & Proponer sugerencias de productos a partir del historial de tratamientos y necesidades detectadas. \\
\hline
\texttt{RF-M5-07} & Salón & Cliente profesional & Gestionar plantillas reutilizables e imágenes del resultado final de los servicios. \\
\hline
\texttt{RF-M5-08} & Salón & Cliente profesional & Generar borradores de correo técnico a partir de la ficha del servicio para contactar con el cliente final. \\
\hline
\RFModuleHeader{M6. Comunicaciones, promociones y formación}
\texttt{RF-M6-01} & Comunicaciones & Administrador, cliente, comercial & Crear, editar, publicar, retirar y consultar promociones comerciales segmentadas por rol, vigencia, producto o cliente, incluyendo adjuntos, descuentos porcentuales, descuentos por volumen y regalos promocionales. \\
\hline
\texttt{RF-M6-02} & Comunicaciones & Administrador, cliente & Gestionar segmentos o categorías jerárquicas de clientes, configurar umbrales de compra para Plata, Oro y Platino, recalcular rangos según compra acumulada del periodo y aplicar promociones por pertenencia heredada, resolviendo como máximo una promoción económica por línea de pedido. \\
\hline
\texttt{RF-M6-03} & Comunicaciones & Administrador, todos los roles & Emitir notificaciones internas, correos, recordatorios y avisos Web Push asociados a eventos relevantes del sistema. \\
\hline
\texttt{RF-M6-04} & Comunicaciones & Administrador, cliente, comercial & Publicar formaciones, mostrar su disponibilidad y permitir la inscripción o cancelación por parte de usuarios autorizados. \\
\hline
\RFModuleHeader{M7. Administración transversal}
\texttt{RF-M7-01} & Administración técnica & Administrador & Agrupar las funciones administrativas reales: usuarios, clientes, catálogo, comunicaciones, pedidos, auditoría, ajustes de avisos automáticos, cargo por agencia y política de umbrales de rangos. \\
\hline
\texttt{RF-M7-02} & Administración técnica & Administrador & Consultar información útil para la gestión diaria sin convertir los diagnósticos técnicos internos en pantallas de uso ordinario. \\
\hline
\texttt{RF-M7-03} & Administración técnica & Administrador & Mantener configuraciones y servicios auxiliares transversales reutilizables por otros módulos. \\
\hline
\texttt{RF-M7-04} & Administración técnica & Administrador & Gestionar parámetros comunes del sistema, incluyendo canales de aviso, importes configurables, umbrales comerciales, estados auxiliares y criterios operativos compartidos. \\
\hline
\texttt{RF-M7-05} & Administración técnica & Administrador & Generar propuestas internas de pedido a proveedor a partir de productos activos, criterios de reposición y datos operativos disponibles en la aplicación. \\
\hline
\texttt{RF-M7-06} & Administración técnica & Administrador & Consultar y exportar auditoría de accesos, acciones administrativas y eventos relevantes en formato trazable. \\
\hline
\texttt{RF-M7-07} & Administración técnica & Equipo técnico & Aplicar y diagnosticar internamente políticas globales de rate limiting sobre acceso y APIs sensibles, reservando su modificación a operación técnica. \\
\hline
\texttt{RF-M7-08} & Administración técnica & Administrador, cliente & Ofrecer soporte transversal de compatibilidad, PWA, navegador no soportado e instalación, manteniendo la comprobación detallada de servicios auxiliares como infraestructura interna. \\
\hline
\RFModuleHeader{M8. Funcionalidades futuras}
\texttt{RF-M8-01} & Funcionalidades futuras & Cliente profesional & Proponer tratamientos mediante un asistente inteligente basado en diagnóstico con capilógrafo, información técnica, historial y contexto del cliente del salón. \\
\hline
\texttt{RF-M8-02} & Funcionalidades futuras & Cliente profesional & Simular estilos o coloraciones mediante la cámara del móvil o imágenes capturadas para apoyar decisiones técnicas futuras. \\
\hline
\texttt{RF-M8-03} & Funcionalidades futuras & Cliente profesional, comercial & Reconocer productos mediante cámara o imagen para consulta, pedido o soporte operativo. \\
\hline
\texttt{RF-M8-04} & Funcionalidades futuras & Cliente final, cliente profesional & Gestionar reservas online de clientes finales vinculadas a salones profesionales. \\
\hline
\texttt{RF-M8-05} & Funcionalidades futuras & Todos los roles & Ofrecer un asistente virtual tipo chatbot para resolver dudas sobre la aplicación, el catálogo, los tratamientos capilares y los procesos comerciales. \\
\hline
\texttt{RF-M8-06} & Funcionalidades futuras & Administrador & Automatizar servicios externos como Factusol mediante flujos con n8n o APIs equivalentes para sincronizar datos y generar documentación comercial futura. \\
\hline
\end{supertabular}
}


\section{Requisitos no funcionales}
\label{sec:NonFunctionalRequirements}
\input{./sections/development/requirements/NonFunctionalRequirements.tex}

\section{Requisitos de información}
\label{sec:InformationRequirements}
\section{Requisitos de información}
\label{sec:InformationRequirements}

En esta sección se resumen los requisitos de información que el sistema debe almacenar, consultar y mantener. Para evitar duplicidad con el modelo conceptual de la \autoref{sec:conceptualmodel}, la Tabla~\ref{tab:information-requirements-executive} presenta el detalle como diccionario compacto: entidad, atributos clave y relaciones principales.

{\scriptsize
\setlength{\parskip}{0pt}
\setlength{\tabcolsep}{3pt}
\renewcommand{\arraystretch}{1.10}
\newcommand{\RIModuleHeader}[1]{%
\hline
\rowcolor{black!8}
\multicolumn{4}{|>{\raggedright\arraybackslash}p{0.93\linewidth}|}{\textbf{#1}}\\
\hline
}
\tablecaption{Diccionario ejecutivo de requisitos de información agrupado por módulos\label{tab:information-requirements-executive}}
\tablefirsthead{%
\hline
\textbf{Módulo} & \textbf{Entidad} & \textbf{Atributos clave} & \textbf{Relaciones principales} \\
\hline}
\tablehead{%
\hline
\textbf{Módulo} & \textbf{Entidad} & \textbf{Atributos clave} & \textbf{Relaciones principales} \\
\hline}
\begin{supertabular}{|>{\raggedright\arraybackslash}p{0.08\linewidth}|>{\raggedright\arraybackslash}p{0.22\linewidth}|>{\raggedright\arraybackslash}p{0.40\linewidth}|>{\raggedright\arraybackslash}p{0.23\linewidth}|}
\RIModuleHeader{M1. Acceso, usuarios y auditoría}
\texttt{RI-M1} & Usuario & Identificador, email, contraseña cifrada, nombre, teléfono, empresa, imagen, último acceso, fechas de creación y actualización. & Rol, estado de usuario, cliente profesional o comercial cuando aplique. \\
\hline
\texttt{RI-M1} & Rol & Código, nombre y descripción del perfil de acceso. & Usuarios asociados y reglas de autorización por interfaz. \\
\hline
\texttt{RI-M1} & Estado de usuario & Código y nombre del estado operativo de la cuenta. & Usuarios y solicitudes administrativas. \\
\hline
\texttt{RI-M1} & Solicitud de registro & Email, datos del solicitante, rol solicitado, origen, estado, fechas de solicitud y revisión, motivo de rechazo. & Estado de solicitud, origen, usuario revisor y usuario creado si se aprueba. \\
\hline
\texttt{RI-M1} & Catálogos de solicitud & Estados y orígenes controlados de las peticiones de alta. & Solicitudes de registro. \\
\hline
\texttt{RI-M1} & Registro de gestión de usuarios & Usuario afectado, administrador, acción, estado/rol anterior y nuevo, motivo, fecha. & Usuario objetivo, usuario administrador y tipo de acción. \\
\hline
\texttt{RI-M1} & Registro de acceso & Usuario, email intentado, tipo de evento, resultado, IP, agente, token de sesión, fecha. & Usuario, tipo de evento y tipo de resultado. \\
\hline
\RIModuleHeader{M2. Clientes profesionales, comerciales y rutas}
\texttt{RI-M2} & Cliente profesional & Usuario vinculado, comercial asignado, nombre, contacto, dirección, ciudad, provincia, geolocalización, ventana de visita y observaciones. & Usuario, comerciales, visitas, pedidos, segmentos y clientes del salón. \\
\hline
\texttt{RI-M2} & Comercial & Usuario vinculado, código, zona, datos operativos, horario y preferencias de ruta. & Usuario, clientes asignados, visitas y rutas. \\
\hline
\texttt{RI-M2} & Asignación comercial-cliente & Cliente, comercial, fecha de asignación, fecha de fin, motivo y usuario responsable. & Cliente profesional y comercial. \\
\hline
\texttt{RI-M2} & Visita comercial & Cliente, comercial, tipo, estado, fecha planificada, resultado, observaciones y contexto de replanificación. & Cliente, comercial, ruta, pedidos de reparto y notificaciones. \\
\hline
\texttt{RI-M2} & Catálogos de visita & Estados y tipos de visita comercial. & Visitas comerciales. \\
\hline
\texttt{RI-M2} & Ruta comercial & Comercial, fecha, estado, origen, destino, distancia, duración estimada y configuración de jornada. & Comercial, estado de ruta y visitas en ruta. \\
\hline
\texttt{RI-M2} & Visita en ruta & Ruta, visita, orden, ETA, distancia acumulada y duración estimada. & Ruta comercial y visita comercial. \\
\hline
\RIModuleHeader{M3. Catálogo, productos y coloración}
\texttt{RI-M3} & Producto & Nombre, referencia, descripción, categoría, línea, subcategoría, estado, imagen e información técnica. & Categoría, línea, subcategoría, estado, recursos, pedidos y productos usados en servicios. \\
\hline
\texttt{RI-M3} & Categoría de producto & Nombre, descripción, imagen y orden de visualización. & Líneas comerciales y productos. \\
\hline
\texttt{RI-M3} & Línea comercial & Categoría, nombre, descripción, imagen y orden. & Categoría, subcategorías, productos y cartas de color. \\
\hline
\texttt{RI-M3} & Subcategoría de producto & Línea, subcategoría padre, nombre, descripción, imagen y orden. & Línea, jerarquía de subcategorías y productos. \\
\hline
\texttt{RI-M3} & Recurso de apoyo & Título, descripción, tipo, URL, producto asociado y orden. & Producto y tipo de recurso. \\
\hline
\texttt{RI-M3} & Carta de color & Nombre, descripción, línea, imagen y orden. & Línea comercial y referencias de color. \\
\hline
\texttt{RI-M3} & Referencia de color & Carta, código, nombre, tono, equivalencia, imagen y orden. & Carta de color, líneas de pedido y fichas técnicas. \\
\hline
\texttt{RI-M3} & Catálogos de catálogo & Estados de producto y tipos de recurso. & Productos y recursos de apoyo. \\
\hline
\RIModuleHeader{M4. Pedidos, reparto y cobro}
\texttt{RI-M4} & Pedido & Cliente, comercial, estado, estado de cobro, importe, método de entrega, cargo de agencia, fechas y notas. & Cliente, comercial, estado, líneas de pedido, repartos y pagos. \\
\hline
\texttt{RI-M4} & Línea de pedido & Pedido, producto, referencia de color, cantidad, precio y observaciones. & Pedido, producto y referencia cromática. \\
\hline
\texttt{RI-M4} & Estados de pedido y cobro & Código, nombre y descripción de estados operativos. & Pedidos. \\
\hline
\texttt{RI-M4} & Reparto de pedido & Pedido, comercial, visita vinculada, estado, método de entrega, bultos manuales, notas, fecha de entrega y usuarios responsables. & Pedido, comercial, visita comercial, líneas de reparto y etiqueta logística. \\
\hline
\texttt{RI-M4} & Línea de reparto & Reparto, línea de pedido y cantidad servida en ese reparto. & Reparto y línea de pedido. \\
\hline
\texttt{RI-M4} & Pago de pedido & Pedido, importe, método, fecha de pago, notas y usuario que registra el cobro. & Pedido y usuario responsable. \\
\hline
\texttt{RI-M4} & Visita de reparto vinculada & Visita comercial que actúa como soporte logístico de entrega para repartos gestionados por comercial. & Reparto, pedido, cliente, comercial y ETA mostrada al cliente. \\
\hline
\RIModuleHeader{M5. Gestión técnica del salón}
\texttt{RI-M5} & Cliente del salón & Cliente profesional propietario, nombre, teléfono, email, observaciones y fechas. & Cliente profesional, servicios, sugerencias y reservas futuras. \\
\hline
\texttt{RI-M5} & Servicio realizado & Cliente del salón, fecha, tipo, descripción, resultado, observaciones e imágenes. & Cliente del salón, ficha técnica, productos usados e imágenes. \\
\hline
\texttt{RI-M5} & Ficha técnica del servicio & Diagnóstico, fórmula, tonalidad, tiempo de exposición, notas técnicas y resultado. & Servicio realizado y referencias de color. \\
\hline
\texttt{RI-M5} & Producto utilizado en servicio & Servicio, producto, referencia cromática, cantidad, uso y observaciones. & Servicio, producto y referencia de color. \\
\hline
\texttt{RI-M5} & Plantilla de servicio & Cliente profesional, nombre, descripción y configuración reutilizable. & Cliente profesional y productos de plantilla. \\
\hline
\texttt{RI-M5} & Imagen de resultado & Servicio, URL, descripción, orden y fecha de subida. & Servicio realizado. \\
\hline
\texttt{RI-M5} & Sugerencia de producto & Cliente del salón, producto, motivo, prioridad y fecha. & Cliente del salón y producto. \\
\hline
\texttt{RI-M5} & Correo técnico generado & Servicio, destinatario, asunto, contenido y fecha de generación. & Servicio y ficha técnica. \\
\hline
\RIModuleHeader{M6. Comunicaciones, promociones y formación}
\texttt{RI-M6} & Promoción & Título, descripción, tipo de descuento, beneficio, porcentaje, umbral de compra, regalo, adjuntos, vigencia, estado, segmento y condiciones. & Productos, segmentos, clientes, tipos de promoción y pedidos con descuento. \\
\hline
\texttt{RI-M6} & Tipo de promoción & Código, nombre, descripción, orden de visualización y estado de actividad. & Promociones comerciales. \\
\hline
\texttt{RI-M6} & Segmento de cliente & Código, nombre, criterio, descripción y fechas. & Clientes profesionales, promociones y asignaciones. \\
\hline
\texttt{RI-M6} & Asignación de segmento & Cliente, segmento, origen, fecha de asignación y usuario responsable. & Cliente profesional y segmento. \\
\hline
\texttt{RI-M6} & Producto promocionado & Promoción, producto y condición de aplicación. & Promoción y producto. \\
\hline
\texttt{RI-M6} & Notificación interna & Destinatario, título, mensaje, tipo, canal, estado de lectura y fechas. & Usuario, cliente, visita, promoción o formación según origen. \\
\hline
\texttt{RI-M6} & Suscripción push & Usuario, endpoint, claves, estado, agente, último uso y revocación. & Usuario y notificaciones Web Push. \\
\hline
\texttt{RI-M6} & Recordatorio & Usuario, título, mensaje, fecha programada, estado y origen. & Usuario y evento funcional relacionado. \\
\hline
\texttt{RI-M6} & Formación e inscripción & Datos de formación, fecha, modalidad, aforo, estado, usuario inscrito y estado de inscripción. & Formación, usuario e integración externa si aplica. \\
\hline
\RIModuleHeader{M7. Administración transversal}
\texttt{RI-M7} & Configuración del sistema & Clave, valor, tipo, ámbito, descripción y fecha de actualización, incluyendo cargo por agencia y política de umbrales/recálculo de rangos. & Usuario administrador y servicios transversales. \\
\hline
\texttt{RI-M7} & Política de rate limiting & Ámbito, límite, ventana temporal, estado y descripción. & Sobrescritura técnica y endpoints protegidos. \\
\hline
\texttt{RI-M7} & Servicio auxiliar transversal & Nombre, tipo, estado, configuración básica, último chequeo y observaciones. & Comunicaciones, configuración del sistema y servicios de soporte. \\
\hline
\texttt{RI-M7} & Registro técnico transversal & Servicio, tipo de operación, estado, resultado, error si existe y fechas. & Servicio auxiliar y eventos técnicos asociados. \\
\hline
\texttt{RI-M7} & Propuesta de pedido a proveedor & Proveedor, estado, criterio de generación, importe estimado, fecha y observaciones. & Productos activos y líneas de propuesta. \\
\hline
\texttt{RI-M7} & Línea de propuesta a proveedor & Propuesta, producto, cantidad, motivo y datos de reposición. & Propuesta y producto. \\
\hline
\RIModuleHeader{M8. Funcionalidades inteligentes previstas}
\texttt{RI-M8} & Diagnóstico con capilógrafo & Usuario, cliente del salón, mediciones o capturas del capilógrafo, tipo y estado del cabello, cuero cabelludo, tratamientos previos, observaciones y validación profesional. & Usuario, cliente del salón y recomendaciones generadas. \\
\hline
\texttt{RI-M8} & Recomendación de tratamiento & Diagnóstico, producto sugerido, descripción, prioridad, motivo, fuente de la recomendación y revisión del profesional. & Diagnóstico capilar y producto. \\
\hline
\texttt{RI-M8} & Simulación de estilo o coloración & Usuario, captura de cámara móvil o imagen origen, cambio solicitado, resultado, producto o tono asociado y fecha. & Usuario y posible referencia de producto. \\
\hline
\texttt{RI-M8} & Reconocimiento de producto & Usuario, imagen, código detectado, producto, propósito y fecha. & Usuario, producto y propósito de reconocimiento. \\
\hline
\texttt{RI-M8} & Reserva online & Cliente profesional, cliente del salón, contacto, servicio, fecha y estado. & Cliente profesional, cliente del salón y estado de reserva. \\
\hline
\texttt{RI-M8} & Conversación con asistente virtual & Usuario, rol, tema, pregunta, respuesta, contexto utilizado, fecha y valoración opcional. & Usuario, catálogo, tratamientos, pedidos o documentación de ayuda según la consulta. \\
\hline
\texttt{RI-M8} & Automatización de servicio externo & Servicio externo, tipo de operación, entidad afectada, payload, resultado, estado, error, reintentos y fecha. & Factusol, n8n, pedidos, clientes, productos y documentos comerciales futuros. \\
\hline
\end{supertabular}
}


\section{Reglas de negocio}
\textcolor{red}{TODO} En el desarrollo del sistema, hay que tener en cuenta las denominadas reglas de negocio, es decir, el conjunto de restricciones, normas o políticas de la organización que deben ser respetadas por el sistema, las cuales suelen ser cambiantes.

\section{Análisis GAP}
\textcolor{red}{TODO} Esta sección es opcional y es de aplicación cuando en la revisión de la técnica se hayan localizado soluciones informáticas que permitan servir de base para el nuevo sistema software. En esta sección se debe identificar y medir las diferencias entre lo que proporcionan dichas soluciones y los requisitos definidos para el proyecto. El resultado de este análisis permitirá identificar cuáles de éstos requisitos ya están solventados total o parcialmente por los distintos sistemas y poder elegir uno de ellos. 